% =====================================================================
%  CARNET D'ENTRAÎNEMENT — Fiche 12 (Chimie organique)
%  THÈME 3 — Nomenclature et priorité des fonctions
%  NIVEAU DIFFICILE — ÉNONCÉS
% =====================================================================
\documentclass[11pt,a4paper]{article}
\usepackage[utf8]{inputenc}
\usepackage[T1]{fontenc}
\usepackage{lmodern}
\usepackage[french]{babel}
\usepackage{amsmath,amssymb,mathtools}
\providecommand{\degree}{^\circ}
\usepackage{icomma}
\usepackage{siunitx}
\sisetup{output-decimal-marker={,},per-mode=symbol}
\usepackage[version=4]{mhchem}
\usepackage{chemfig}
\setchemfig{atom sep=2em,bond length=0.9em,cram width=2.5pt}
\usepackage{xcolor}
\usepackage{array}
\usepackage{booktabs}
\usepackage{tabularx}
\usepackage[most]{tcolorbox}
\usepackage{enumitem}
\usepackage{tikz}
\usetikzlibrary{arrows.meta,positioning,calc}
\usepackage{fancyhdr}
\usepackage{titlesec}
\usepackage[hidelinks]{hyperref}
\usepackage[a4paper,margin=2.1cm,headheight=26pt,headsep=9pt]{geometry}

\definecolor{primary}{RGB}{150,98,162}
\definecolor{primarydark}{RGB}{92,52,110}
\definecolor{excol}{RGB}{0,135,90}
\definecolor{attncol}{RGB}{200,40,55}
\definecolor{methcol}{RGB}{90,90,100}

\newcommand{\runtitle}{Thème 3 — Difficile — Énoncés}
\pagestyle{fancy}\fancyhf{}
\fancyhead[L]{\small\textcolor{primarydark}{\textbf{Fiche 12} — Chimie organique}}
\fancyhead[R]{\small\textcolor{primarydark}{\runtitle}}
\fancyfoot[C]{\small\thepage}
\renewcommand{\headrulewidth}{0.8pt}
\renewcommand{\headrule}{\hbox to\headwidth{\color{primary}\leaders\hrule height \headrulewidth\hfill}}
\setlength{\parindent}{0pt}\setlength{\parskip}{3pt}

\tcbset{boxbase/.style={enhanced,breakable,boxrule=0.9pt,arc=2.5pt,
   left=3mm,right=3mm,top=2mm,bottom=2mm,fonttitle=\bfseries,coltitle=white,
   toptitle=0.5mm,bottomtitle=0.5mm}}
\newtcolorbox[auto counter]{exo}[1]{boxbase,colback=primary!4,colframe=primary,
   title={\textsc{Exercice}~\thetcbcounter\ \textmd{--- \itshape #1}}}
\newtcolorbox{consigne}{boxbase,colback=methcol!8,colframe=methcol,sharp corners,
   title={\textsc{Consignes}}}

\newcommand{\banner}[2]{%
\begin{tcolorbox}[enhanced,colback=primary,colframe=primary,arc=5pt,boxrule=0pt,
   left=5mm,right=5mm,top=2.5mm,bottom=2.5mm,coltext=white]
{\large\bfseries #1}\\[1pt]{\small #2}
\end{tcolorbox}\vspace{2pt}}

\begin{document}

\banner{Thème 3 — Nomenclature et priorité des fonctions}{Niveau \textbf{Difficile} — Énoncés \quad$\vert$\quad Fiche 12, Chimie organique}

\begin{consigne}
Exercices \textbf{ouverts} à \textbf{sous-questions progressives}. Applique la méthode complète et justifie chaque étape (priorité, choix de la chaîne, sens de numérotation).
\end{consigne}

% ---------------------------------------------------------------------
\begin{exo}{nomenclature complète d'une molécule trifonctionnelle}
On considère la molécule suivante :
\begin{center}\chemfig{HO-[:30](=[:90]O)-[:-30]CH_2-[:30]CH(-[:90]OH)-[:-30]CH(-[:-90]NH_2)-[:30]CH_3}\end{center}
\begin{enumerate}[label=\textbf{\alph*)},leftmargin=2em,itemsep=1pt]
  \item Repère les \textbf{trois} fonctions présentes.
  \item Classe-les par priorité : quelle est la fonction principale (suffixe) ? Quelles fonctions deviennent des préfixes, et lesquels ?
  \item Choisis la chaîne principale et numérote-la (justifie le n\textsuperscript{o}1).
  \item Donne le \textbf{nom complet} (préfixes par ordre alphabétique).
\end{enumerate}
\end{exo}

% ---------------------------------------------------------------------
\begin{exo}{du nom à la structure (sens inverse)}
On donne le nom \textbf{acide 2-aminopropanoïque} (c'est l'\emph{alanine}, un acide aminé).
\begin{enumerate}[label=\textbf{\alph*)},leftmargin=2em,itemsep=1pt]
  \item Décode le nom (chaîne, position du groupe amino, fonction principale) et \textbf{dessine} la structure semi-développée.
  \item Identifie les deux fonctions et vérifie la cohérence avec la priorité ($\ce{-COOH}$ principale, $\ce{-NH2}$ en préfixe amino-).
\end{enumerate}
\end{exo}

% ---------------------------------------------------------------------
\begin{exo}{alcool + double liaison : locants et suffixe -èn-ol}
Soit la molécule :
\begin{center}\chemfig{CH_2=[:30]CH-[:-30]CH_2-[:30]CH(-[:90]OH)-[:-30]CH_3}\end{center}
\begin{enumerate}[label=\textbf{\alph*)},leftmargin=2em,itemsep=1pt]
  \item Identifie les deux fonctions et indique laquelle est prioritaire.
  \item Numérote la chaîne dans les deux sens ; donne à chaque fois la position de l'$\ce{-OH}$ et celle de la double liaison.
  \item Quel sens retient-on ? Donne le nom (la double liaison se marque par \textit{-én-} avec son indice, l'alcool par \textit{-ol}).
\end{enumerate}
\end{exo}

% ---------------------------------------------------------------------
\begin{exo}{nommer un ester}
On donne la molécule ci-dessous :
\begin{center}\chemfig{CH_3-[:30]CH_2-[:-30]CH_2-[:30](=[:90]O)-[:-30]O-[:30]CH_3}\end{center}
\begin{enumerate}[label=\textbf{\alph*)},leftmargin=2em,itemsep=1pt]
  \item Quelle fonction reconnais-tu ? Justifie par le voisinage du carbonyle.
  \item Un ester se nomme «\,\emph{(acide)}-oate de \emph{(alkyl)}-yle\,». Identifie la partie issue de l'acide et la partie issue de l'alcool.
  \item Donne le nom complet de l'ester.
\end{enumerate}
\end{exo}

\end{document}
